\section{Several competing templates and a mass threshold}
\label{sec:multi_competitor}
This result extends the fixed-state three-template comparison.
All components have the common isotropic covariance $\tau^2I$.
Let $\mu_0=0$, $\mu_1=e$ with $\norm e=1$, and fix $x$ with
$s=e^\top x\in(0,1/2)$ and $P=1/2-s$.
For $M\ge1$ competitors $\nu_j$, suppose
\begin{align}
 e^\top\nu_j&=-a_j,\qquad a_j>0,\nonumber\\
 Q_j&=\tfrac12(\norm{x-\nu_j}^2-\norm x^2)>P.
 \label{eq:multi_geometry}
\end{align}
Their prior ratios are $1:r:hq_1:\cdots:hq_M$ where
$r,h,q_j>0$ and $\sum_jq_j=1$.
The internal composition $q$ remains fixed as the total ratio $h$
varies. Define
\begin{align}
 \bar a&=\sum_jq_ja_j,\quad
 \bar Q=\sum_jq_jQ_j,\quad
 \overline{aQ}=\sum_jq_ja_jQ_j,\nonumber\\
 C_s&=(1-s)r-s>0,\nonumber\\
 h_{\max}&=\frac{C_s}{\bar a+s},\qquad h_0=\frac r{\bar a}.
 \label{eq:multi_mass_range}
\end{align}
Let $z_L=1/v_L$ with
\begin{equation}
 v_L>\tau^2,\qquad r e^{-Pz_L}\le \frac{s}{1-s}.
 \label{eq:multi_endpoint_conditions}
\end{equation}
For inverse variance $z$, put $l(z)=1-\tau^2z$ and
\begin{align}
 m_h(z)&=
 \frac{r e^{-Pz}-h\sum_j a_jq_je^{-Q_jz}}
      {1+r e^{-Pz}+h\sum_jq_je^{-Q_jz}},\nonumber\\
 g_h(z)&=l(z)(m_h(z)-s).
 \label{eq:multi_drift}
\end{align}
The projected RED displacement per unit step size is
$e^\top(b-Bx)+\lambda g_h(z)$.
Thus comparison of $g_h$ compares actual updates at this state,
for the same $\lambda>0$, without alignment of $e$ and $B$.

\begin{proposition}[Mass threshold with several competitors]
\label{prop:multi_mass}
Under \eqref{eq:multi_geometry}--\eqref{eq:multi_endpoint_conditions},
for $0<h\le h_{\max}$, the infinite-variance endpoint
$g_h(0)$ dominates the terminal endpoint $g_h(z_L)$.
A finite intermediate variance strictly improves on both endpoints
if and only if $D(h)>0$, where
\begin{align}
 D(h)={}&-Pr\nonumber\\
 &+h[(1+r)\overline{aQ}+r\bar Q-Pr(1+\bar a)]\nonumber\\
 &+h^2(\overline{aQ}-\bar a\bar Q)\nonumber\\
 &-\tau^2[C_s-h(\bar a+s)](1+r+h).
 \label{eq:multi_polynomial}
\end{align}
There is exactly one root $h_{\rm crit}\in(0,h_0)$.
On $0<h\le h_{\max}$, improvement is equivalent to
$h>h_{\rm crit}$. At fixed remaining scalar inputs, increasing
any $Q_j$ strictly decreases $h_{\rm crit}$.
\end{proposition}

\begin{proof}
The mass range gives
\[
 u:=g_h(0)=\frac{C_s-h(\bar a+s)}{1+r+h}\ge0.
\]
Without the competitors, the terminal right-template responsibility
is at most $s$ by \eqref{eq:multi_endpoint_conditions}.
Adding negative-projection templates decreases its horizontal mean,
so $g_h(z_L)\le0$ and the high-variance endpoint is larger.

Set
\begin{align*}
 A_u(z)&=(1-s)l(z)-u,\\
 U_u(z)&=\frac{s l(z)+u}{A_u(z)},\\
 W_{u,j}(z)&=\frac{(a_j+s)l(z)+u}{A_u(z)}.
\end{align*}
Since $u\le C_s/(1+r)<1-s$, $A_u(0)>0$.
The set $I_u=\{z\in[0,z_L]:A_u(z)>0\}$ is an initial interval.
On that interval, multiplication by the positive responsibility
denominator gives the equivalence
\begin{align}
 g_h(z)>u
 &\quad\Longleftrightarrow\quad J_u(z)<r,\nonumber\\
 J_u(z)&=U_u(z)e^{Pz}
     +h\sum_jq_jW_{u,j}(z)e^{-(Q_j-P)z}.
 \label{eq:multi_convex_test}
\end{align}
In particular $J_u(0)=r$.
The identities
\begin{align*}
 U_u&=\frac{s}{1-s}+\frac{u}{(1-s)A_u},\\
 W_{u,j}&=\frac{a_j+s}{1-s}
             +\frac{(1+a_j)u}{(1-s)A_u}
\end{align*}
show that each ratio is log-convex.
Indeed a positive affine reciprocal is log-convex on its positive
domain, and positive sums preserve log-convexity.
Multiplication by an exponential preserves that property.
Moreover $U_u>0$, $U_u',U_u''\ge0$ and $P>0$, so the first
term in $J_u$ is strictly convex. Thus $J_u$ is strictly convex.

No improvement is possible where $A_u(z)\le0$.
All template projections are at most 1 and the weight on $\mu_0$
is positive, giving $m_h(z)<1$. Since $l(z)>0$,
\[
 g_h(z)-u<(1-s)l(z)-u=A_u(z)\le0.
\]
Consequently an improvement anywhere in $(0,z_L)$ exists exactly
when $J_u'(0)<0$. This argument does not require $A_u$ to remain
positive all the way to $z_L$.

On $I_u$ the exact identity
\[
 g_h(z)-u=
 \frac{A_u(z)e^{-Pz}[r-J_u(z)]}
      {1+r e^{-Pz}+h\sum_jq_je^{-Q_jz}}
\]
shows that the derivatives of $g_h$ and $J_u$ at zero have
opposite signs. Direct differentiation of
\eqref{eq:multi_drift} yields
\[
 g_h'(0)=\frac{D(h)}{(1+r+h)^2},
\]
which proves the improvement criterion.

For uniqueness,
$D(0)=-Pr-\tau^2C_s(1+r)<0$.
At $h_0=r/\bar a$, $m_{h_0}(0)=0$ and
\[
 m_{h_0}'(0)=
 \frac{r(\overline{aQ}/\bar a-P)}{1+r+h_0}>0.
\]
Hence $g_{h_0}'(0)=m_{h_0}'(0)+\tau^2s>0$ and $D(h_0)>0$.
A polynomial of degree at most two with opposite signs at these
endpoints has exactly one root inside, which is simple and crosses
from negative to positive. Since $h_{\max}<h_0$, this proves the
claimed threshold on the comparison range.
Finally,
\[
 \frac{\partial D(h)}{\partial Q_j}
 =h q_j[a_j(1+r+h)+(r-h\bar a)]>0
 \quad(0<h<h_0).
\]
At the root $\partial_hD>0$, so implicit differentiation gives
$\partial_{Q_j}h_{\rm crit}<0$.
\end{proof}

The comparison of a partial derivative in $Q_j$ keeps $a_j$
and the other scalar inputs fixed. A geometric perturbation can
change several of these inputs at once.
For $M=1$, the covariance term in \eqref{eq:multi_polynomial}
vanishes and the three-template polynomial is recovered.
For $M>1$, $\operatorname{Cov}_q(a,Q)$ can be negative, so
global uniqueness on the positive $h$ axis is not asserted.

\paragraph{A realizable second positive root.}
Take $x=(.3,3)$, $e=(1,0)$, $r=1.29$,
$\tau^2=10^{-4}$, $v_L=.01$ and $q_1=q_2=1/2$, with
\begin{align*}
 \nu_1&=(-.1,3+\sqrt{28.93}),\\
 \nu_2&=(-3,3+\sqrt{2.2}).
\end{align*}
Then $a=(.1,3)$ and $Q=(10,2)$ give
\[
 D(h)=-5.799815h^2+15.09746335h-.258138087.
\]
Its positive roots are approximately $.01721192$ and $2.58588204$,
while $h_{\max}\simeq.32594595<h_0\simeq.83225806$.
Only the first root belongs to the threshold interval.

\paragraph{Scope and numerical checks.}
The result is a fixed-state single-step comparison.
It does not provide a channel, a terminal entrance, or a raw-output
risk separation for arbitrary competitor configurations.
Strict improvement over the infinite-variance endpoint persists
against a sufficiently large finite high variance by continuity.
The program \path{experiments/code/theory_extension_checks.py}
checks the derivative identity and convex criterion on 128 finite
cases, including 69 cases outside the former sufficient condition
$C_s/(1+r)<(1-s)l(z_L)$.
It also checks the explicit two-root geometry, the single-competitor
reduction and 64 four-corner attribution identities.
These checks are distinct from the external draft's reported
1500-case batch and do not replace the analytic proof.
